\begin{table}[htbp]
\centering
\caption{Tail-estimate comparison for VaR$_{99.5\%}$ of the transferred-severity distributions: empirical, frequentist extreme-value (peaks-over-threshold, POT) and Bayesian POT. Point estimate with 95\% interval; POT threshold at the 90th percentile ($N_u\approx69$ exceedances).}
\label{tab:tail_comparison}
\begin{tabular}{lcc}
\toprule
\textbf{Method} & \textbf{V1 (adjusted)} & \textbf{V2 (new profile)} \\
\midrule
Empirical & 0.314 [0.227, 0.432] & 0.302 [0.221, 0.411] \\
EVT, frequentist POT & 0.305 [0.229, 0.368] & 0.293 [0.220, 0.355] \\
EVT, Bayesian POT & 0.314 [0.267, 0.419] & 0.302 [0.257, 0.403] \\
\bottomrule
\end{tabular}
\vspace{2pt}
{\footnotesize The three estimates are mutually consistent: each point lies inside the other methods' intervals, and both EVT intervals contain the empirical point. The sign of the fitted tail shape is not resolved ($\hat\xi\approx0.07$ for V1, $0.07$ for V2; 95\% credible intervals $[-0.16,+0.38]$ and $[-0.16,+0.38]$, both spanning zero), and the upper limits are wide. Each row is a different inferential object. The empirical row is a credible interval under the donor-composition posterior (bayesian bootstrap by syndicate x posterior draws, Dirichlet(1, ..., 1) over syndicates (Rubin's, unmodified)), with one posterior draw per replicate. The frequentist-POT row is conditional frequentist: donor resampling only, theta fixed at its posterior mean. The Bayesian-POT row is the GPD posterior.}
\end{table}
